\documentclass[12pt,a4paper, oneside]{article}
\usepackage{graphicx,amssymb,amsfonts,latexsym,amsmath,amsthm,times}
\usepackage{epsfig}
\usepackage{fancyhdr}
\usepackage{color}
%\usepackage[english]{babel} % usar em artigos escritos em Inglês
\usepackage[portuguese]{babel}
\usepackage{authblk}
\usepackage{fancyhdr}
\usepackage[utf8]{inputenc}
\usepackage[T1]{fontenc}
\usepackage{cite}
\usepackage{url}

%Cabeçalho
\fancyhead{}%limpar o head
\fancyhead[L]{\em \textbf{ESTG Masters} -- 7.Março.2018 -- CIICESI, ESTG / P.PORTO}



\renewcommand{\headrulewidth}{0pt}
\renewcommand{\footrulewidth}{0cm}
\setlength\headheight{3cm}

\pagestyle{fancy}


\addtolength{\voffset}{-2.0cm} \addtolength{\textheight}{0.0cm} 

	
\renewcommand\Authfont{\scshape}


\newcommand{\smalllineskip}{\baselineskip=15pt}
\newcommand{\keywords}[1]{{\footnotesize\hspace{0.68cm}{\textit{Keywords}: }#1\par
  \vskip.7\baselineskip}}
\renewenvironment{abstract}[0]{
        \begin{center}\rm\bf{RESUMO}\large\rm
        \\ \vspace{8pt}
        \begin{minipage}{5.2in}\smalllineskip 
        \hspace{1pc}}{\end{minipage}\end{center}\vspace{-1pt}}
\newcommand{\emailaddress}[1]{\newline{\sf#1}}



\setlength{\textwidth}{6.5in} \textheight=8.5in \oddsidemargin 0in \topmargin 0.3 in




\begin{document}

\title{Template para preparação de Artigos do ESTG Masters}
\author{Nome do Estudante$^{1,2}$, Nome do orientador$^2$, Nome do co-orientador$^3$}
\date{%
\small
    $^1$Mestrado em XXXX\\
    $^2$CIICESI, ESTG /P.PORTO -- 
    Centro de Inovação e Investigação em Ciências Empresariais e Sistemas de Informação, Escola Superior de Tecnologia e Gestão / Politécnico do Porto\\
    	$^3$ Outra, se existir\\
emails: 80000@estg.ipp.pt, orientador@estg.ipp.pt, co@estg.ipp.pt
}



\maketitle

\begin{abstract}
Escrever aqui o resumo (150 palavras) \ldots

\textbf{Palavras-chave}: Escrever aqui as palavras-chave (no máximo 6).
\end{abstract}


\section{Introdução}\label{sec:int}

Escrever texto\ldots

Para citar a secção pode escrever-se \verb|Secção \ref{int}| 
e obtém-se a referência automática: Secção \ref{sec:int}.

\subsection{Apresentação e Oportunidade do Tema}\label{subsec:ap}

\subsection{Objetivos principais}\label{subsec:obj}

\subsection{Contributos inovadores}\label{subsec:conino}

\subsubsection{Introduzir figuras}
\begin{figure}[htbp]
\centering
\includegraphics[width=0.8\textwidth]{ESTGmasters}
\caption{ESTGmasters logo}
\label{fig:logo}
\end{figure}

Para citar a Figura \ref{fig:logo} .

\subsubsection{Introduzir Tabelas}

Pode introduzir tabelas usando a página \url{http://www.tablesgenerator.com/#} ou use o exemplo da Tabela \ref{my-label}.

\begin{table}[]
\centering
\caption{My caption}
\label{my-label}
\begin{tabular}{c|cccc}
  & \textbf{1} & \textbf{2} & \textbf{3} & \textbf{4} \\ \hline
A & \multicolumn{3}{c}{branco}           & preto      \\
B & amarelo    & \multicolumn{3}{c}{branco}           \\
C & azul       & preto      & branco     & vermelho   \\ \hline
\end{tabular}
\end{table}

\section{Revisão da literatura}\label{sec:rev}
Para citar Blah blah blah blah blah blah by Bob ~\cite{BobsBook01}.

\subsection{Temáticas teóricas da dissertação} \label{subsec:con}

(ex. grandes temas da tese )

\subsection{Metodologias da Investigação}\label{sec:met}



\section{Conclusões}\label{sec:con}
We worked hard, and achieved very little.

\bibliography{mybib}{}
\bibliographystyle{plain} % Referências estilo habitual
%\bibliographystyle{apalike} % Referências estilo APA

\end{document}

Se estiver texto aqui, depois de \end{document} não será impresso....
